\documentclass[xetex,mathserif,serif,handout]{beamer}

\usepackage{xunicode}
\usepackage{xltxtra}
\usepackage{color}
\usepackage{url}
\usepackage{listings}
\usepackage{fontspec}
\usepackage{geometry}
\usepackage{lastpage}
\usepackage{fancyhdr}
\usepackage{amsmath}
\usepackage{amsthm}
\usepackage{amssymb}
\usepackage{blkarray}
\usepackage{multicol}
\usepackage{relsize}

\definecolor{solarized@base03}{HTML}{002B36}
\definecolor{solarized@base02}{HTML}{073642}
\definecolor{solarized@base01}{HTML}{586e75}
\definecolor{solarized@base00}{HTML}{657b83}
\definecolor{solarized@base0}{HTML}{839496}
\definecolor{solarized@base1}{HTML}{93a1a1}
\definecolor{solarized@base2}{HTML}{EEE8D5}
\definecolor{solarized@base3}{HTML}{FDF6E3}
\definecolor{solarized@yellow}{HTML}{B58900}
\definecolor{solarized@orange}{HTML}{CB4B16}
\definecolor{solarized@red}{HTML}{DC322F}
\definecolor{solarized@magenta}{HTML}{D33682}
\definecolor{solarized@violet}{HTML}{6C71C4}
\definecolor{solarized@blue}{HTML}{268BD2}
\definecolor{solarized@cyan}{HTML}{2AA198}
\definecolor{solarized@green}{HTML}{859900}
\definecolor{yaleblue}{HTML}{0E4C92}

\setbeamertemplate{navigation symbols}{}
% \setbeamerfont{title}{family=\old}
% \setbeamerfont{author}{family=\tfont}%
% \setbeamerfont{frametitle}{family=\oldA}
% \setbeamerfont{date}{family=\dfont}

\setbeamertemplate{itemize items}{--}
\setbeamercolor*{item}{fg=black}

\defaultfontfeatures{Mapping=tex-text}
\hypersetup{pdfstartview={FitH}}

\newcommand{\old}[1]{\fontspec[Alternate=1,Ligatures={Common}]{Hoefler Text}\fontsize{18pt}{30pt}\selectfont #1}%
\newcommand{\oldA}[1]{\fontspec[Alternate=1,Ligatures={Common, Rare}]{Hoefler Text}\fontsize{12pt}{15pt}\selectfont #1}%
\newcommand{\oldB}[1]{\fontspec[Ligatures={Common}]{Didot}\fontsize{12pt}{15pt}\color{solarized@base02}\selectfont #1}%
\newcommand{\tfont}[1]{\fontspec[Alternate=1,Ligatures={Common}]{Hoefler Text}\fontsize{12pt}{20pt}\selectfont #1}%
\newcommand{\dfont}[1]{\fontspec[Ligatures={Common}]{Didot}\fontsize{12pt}{12pt}\selectfont #1}%

\newcommand{\minimize}{\mathop{\mathrm{minimize}}}
\newcommand{\argmin}{\mathop{\mathrm{arg\,min}}}
\newcommand{\argmax}{\mathop{\mathrm{arg\,max}}}
\newcommand{\st}{\mathop{\mathrm{subject\,\,to}}}

\newcommand\independent{\protect\mathpalette{\protect\independenT}{\perp}}
\def\independenT#1#2{\mathrel{\rlap{$#1#2$}\mkern2mu{#1#2}}}

\setlength{\parindent}{0pt}
\setlength{\parskip}{12pt}

\setromanfont [Ligatures={Common}, Numbers={OldStyle}, Variant=01,
 BoldFont={LinLibertine_RB.otf},
 ItalicFont={LinLibertine_RI.otf},
 BoldItalicFont={LinLibertine_RBI.otf}
 ]{LinLibertine_R.otf}



\begin{document}

%%%%%%%%%%%%%%%%%%%%%%%%%%%%%%%%%%%%%%%%%%%%%%%%%%%
\begin{frame}[fragile] \frametitle{}

\vfill

{\fontsize{0.7cm}{0cm}\selectfont Lecture 09 \\\vspace{0.2cm}
Prediction and Leverage with ASA Flight Data}\\\vspace{0.5cm}
30 September 2015

\vspace{2cm}

\begin{minipage}{0.6\textwidth}
Taylor B. Arnold \\
Yale Statistics \\
STAT 312/612
\end{minipage}
\hfill
\begin{minipage}{0.3\textwidth}\raggedleft
\includegraphics[scale=0.3]{../yale-logo.png}
\end{minipage}%

\end{frame}


%%%%%%%%%%%%%%%%%%%%%%%%%%%%%%%%%%%%%%%%%%%%%%%%%%%
\begin{frame}[fragile] \frametitle{}

{\color{yaleblue}\fontsize{16pt}{20pt}\selectfont Goals for today}

\begin{itemize}
\item Quantify leverage of an observation
\item Formulate prediction and confidence intervals for multivariate regression
\item Apply to ASA airline dataset
\end{itemize}

\end{frame}

%%%%%%%%%%%%%%%%%%%%%%%%%%%%%%%%%%%%%%%%%%%%%%%%%%%
\begin{frame}[fragile] \frametitle{}

\begin{flushright}
{\color{yaleblue}\sc\fontsize{1cm}{0cm}\selectfont Leverage}
\end{flushright}

\end{frame}

%%%%%%%%%%%%%%%%%%%%%%%%%%%%%%%%%%%%%%%%%%%%%%%%%%%
\begin{frame}[fragile] \frametitle{}

Recall that {\it leverage} was generally defined as the
amount of influence a point has the estimation of
$\widehat{\beta}$.

\pause We can formally define leverage as the diagonal
elements of the projection matrix:
\begin{align*}
l_i &= P_{ii} \\
&= \left[ X (X^t X)^{-1} X^t \right]_{ii}
\end{align*}

\end{frame}

%%%%%%%%%%%%%%%%%%%%%%%%%%%%%%%%%%%%%%%%%%%%%%%%%%%
\begin{frame}[fragile] \frametitle{}

Notice that $l_{i}$ will be a number between $0$ and $1$;
because $P$ is idempotent and symmetric:
\begin{eqnarray*}
l_{i} &=& \sum_j p_{i,j} p_{j,i} \\ \pause
&=& \sum_j p_{i,j}^2 \\ \pause
&=& p_{i,i}^2 + \sum_{j \neq i} p_{i,j}^2 \\ \pause
&=& l_{i}^2 + \sum_{j \neq i} p_{i,j}^2
\end{eqnarray*}
\pause So then $l_i \geq l_i^2$, which shows the bounds on the
leverage values.

\end{frame}

%%%%%%%%%%%%%%%%%%%%%%%%%%%%%%%%%%%%%%%%%%%%%%%%%%%
\begin{frame}[fragile] \frametitle{}

To see why we use the diagonal of the projection matrix,
look at the variance of the residuals:
\begin{eqnarray*}
\mathbb{V} (r | X) &=& \mathbb{E} (r r^t | X) \\ \pause
&=& \mathbb{E} (M\epsilon \epsilon^t M^t | X) \\ \pause
&=& M \cdot \mathbb{E} (\epsilon \epsilon^t| X) \cdot  M^t \\ \pause
&=& \sigma^2 M M^t \\ \pause
&=& \sigma^2 \left[ \mathbb{I}_n - P \right]
\end{eqnarray*}
\pause So the variance of an individual residual is $\sigma^2(1 - l_i)$, so
for a leverage close to $1$ the regression line will generally pass
very close to the point $i$.

\end{frame}

%%%%%%%%%%%%%%%%%%%%%%%%%%%%%%%%%%%%%%%%%%%%%%%%%%%
\begin{frame}[fragile] \frametitle{}

The individual variance of the $i$'th residual suggests that we
could standardize each residual as such:
\begin{align*}
\frac{r_i}{\sqrt{s^2 (1 - l_i)}}
\end{align*}
\pause This is known as the Studentized residual. If $s^2$ is modifed
to be calculated without the point $i$, externally studentized, then
this quantity follows a $t$ distribution with $n-p$ degrees of freedom.
(Note: This is actually very easy to prove given our already established
results.)

\end{frame}

%%%%%%%%%%%%%%%%%%%%%%%%%%%%%%%%%%%%%%%%%%%%%%%%%%%
\begin{frame}[fragile] \frametitle{}

\begin{flushright}
{\color{yaleblue}\sc\fontsize{1cm}{0cm}\selectfont Confidence and Prediction Intervals}
\end{flushright}

\end{frame}

%%%%%%%%%%%%%%%%%%%%%%%%%%%%%%%%%%%%%%%%%%%%%%%%%%%
\begin{frame}[fragile] \frametitle{}

Now, we consider the case where we observe a new set of
observations that were not used in the estimation of $\widehat{\beta}$:
\begin{align*}
y_{new} &= X_{new} \beta + \epsilon_{new}
\end{align*}
Often we do not actually observe the new values of $y$, but wish
to estimate them from the estimate of $\beta$ and the new
data points.

\end{frame}

%%%%%%%%%%%%%%%%%%%%%%%%%%%%%%%%%%%%%%%%%%%%%%%%%%%
\begin{frame}[fragile] \frametitle{}

An obvious estimate, $\hat{y}_{new}$, is $X_{new} \widehat{\beta}$.
\pause We see easily that:
\begin{align*}
\mathbb{E} (\hat{y}_{new} | X) &= \mathbb{E} (X_{new} \widehat{\beta} | X) \\
&= X_{new} \beta
\end{align*}
Where the conditional on $X$ is with respect to the original data and
the new data matrix $X_{new}$.

\end{frame}

%%%%%%%%%%%%%%%%%%%%%%%%%%%%%%%%%%%%%%%%%%%%%%%%%%%
\begin{frame}[fragile] \frametitle{}

What would we do if we needed a confidence interval for where
the values $y_{new}$ should be located? \pause We need to
calculate the variance of our estimator:
\begin{align*}
\mathbb{V} (\hat{y}_{new} | X) &= \mathbb{V} ( X_{new} \widehat{\beta} | X ) \\
&= X_{new} \mathbb{V} (\widehat{\beta} | X) X_{new}^t \\
&= \sigma^2 X_{new} (X^t X)^{-1} X_{new}^t \\
\end{align*}

\end{frame}

%%%%%%%%%%%%%%%%%%%%%%%%%%%%%%%%%%%%%%%%%%%%%%%%%%%
\begin{frame}[fragile] \frametitle{}

Notice that in the special case that row $j$ of $X_{new}$ is equal
to row $i$ of $X$, we have:
\begin{align*}
\mathbb{V} (\left[ \hat{y}_{new} \right]_j | X) &= \sigma^2 P_{ii} \\
&= \sigma^2 l_i
\end{align*}
So points with high leverage are points where predictions are
particularly variable. \pause Counterintuitive?

\end{frame}

%%%%%%%%%%%%%%%%%%%%%%%%%%%%%%%%%%%%%%%%%%%%%%%%%%%
\begin{frame}[fragile] \frametitle{}

This suggests that we construct the following confidence interval
for the mean of $y_{new}$:
\begin{align*}
\widehat{\mathbb{E}(y_{new} | X)} &\in X_{new} \widehat{\beta} \pm t_{n-p,1-\alpha/2} \cdot
\sqrt{s^2 X_{new} (X^t X)^{-1} X_{new}^t} \\
\end{align*}

\end{frame}

%%%%%%%%%%%%%%%%%%%%%%%%%%%%%%%%%%%%%%%%%%%%%%%%%%%
\begin{frame}[fragile] \frametitle{}

Typically, we are interested in an interval for the actually
observations $y_{new}$ rather than the mean of $y_{new}$.

\pause To calculate the variance of the actual prediction,
see that ($k$ is the number of row of $X_{new}$)
\begin{align*}
\mathbb{V} (y_{new} - \hat{y}_{new} | X) &= \mathbb{V} (y_{new} | X) + \mathbb{V} (\hat{y}_{new} | X) \\
&= \sigma^2 I_k + \sigma^2 X_{new} (X^t X)^{-1} X_{new}^t \\
&= \sigma^2 \left[I_k + X_{new} (X^t X)^{-1} X_{new}^t \right]
\end{align*}

\end{frame}

%%%%%%%%%%%%%%%%%%%%%%%%%%%%%%%%%%%%%%%%%%%%%%%%%%%
\begin{frame}[fragile] \frametitle{}

From here, we now have the following prediction interval:
\begin{align*}
y_{new} | X &\in X_{new} \widehat{\beta} \pm t_{n-p,1-\alpha/2} \cdot
\sqrt{s^2 \left[ I_k + X_{new} (X^t X)^{-1} X_{new}^t \right]} \\
\end{align*}
Which is exactly a factor of $s$ wider than the confidence interval.


\end{frame}

%%%%%%%%%%%%%%%%%%%%%%%%%%%%%%%%%%%%%%%%%%%%%%%%%%%
\begin{frame}[fragile] \frametitle{}

\begin{flushright}
{\color{yaleblue}\sc\fontsize{1cm}{0cm}\selectfont Application to ASA Data}
\end{flushright}

\end{frame}

\end{document}











